\section{Bonding, Lewis Structures, Formal Charge, and VSEPR}

\subsection{Ctrl+F keywords: Lewis, formal charge, resonance, VSEPR, bond angle, planar, polarity, dipole, hybridization}
Use for drawing or checking Lewis structures, choosing resonance forms, molecular geometry, bond angles, polarity, and planar/non-planar questions.

\subsection{Lewis structures}

\subsubsection{Valence electron count}
\[
N_{\mathrm{val}}=\sum_i N_{\mathrm{val},i}-q
\]
For an ion, subtract positive charge and add negative charge. Total Lewis electrons must equal \(N_{\mathrm{val}}\).

\subsubsection{Formal charge}
\[
\mathrm{FC}=V-\left(N_{\mathrm{nonbonding}}+\frac12N_{\mathrm{bonding}}\right)
\]
Preferred structures usually minimize formal charges and place negative formal charge on more electronegative atoms.

\subsubsection{Resonance structures}
\[
\text{resonance forms differ in electron placement, not atom placement}
\]
Equivalent resonance forms indicate delocalization and equalized bond lengths. Do not move nuclei when drawing resonance.

\subsubsection{Octet, expanded octet, electron-deficient exceptions}
\[
\text{period 2 atoms normally obey octet}, \qquad
\text{period }3+\text{ may expand octet}
\]
Hydrogen has a duet. Boron can be electron deficient; sulfur, phosphorus, iodine, and xenon can appear hypervalent in course examples.

\subsection{Exam recipe: Lewis structure and formal charge}
\textbf{Use when:} the problem asks for nitrate, oxalate, resonance, formal charge, or a valid Lewis structure.

\textbf{Steps:}
\begin{enumerate}
    \item Count total valence electrons, including ion charge.
    \item Choose the least electronegative central atom, never hydrogen.
    \item Connect atoms with single bonds and complete outer octets.
    \item Put leftover electrons on the central atom if allowed.
    \item If the central atom lacks an octet, form multiple bonds.
    \item Calculate formal charges and choose the best resonance form.
\end{enumerate}

\textbf{Fast checks:} electron count matches; total formal charge equals ion charge; second-period atoms do not exceed octet.

\subsection{VSEPR geometry}

\subsubsection{Steric number and electron domains}
\[
\mathrm{SN}=\#\text{bonded atoms}+\#\text{lone pairs on central atom}
\]
Multiple bonds count as one electron domain in VSEPR. Lone pairs compress bond angles.

\subsubsection{Common VSEPR geometries and bond angles}
\[
\begin{array}{c c c c}
\toprule
\mathrm{SN} & \text{lone pairs} & \text{molecular geometry} & \text{angle}\\
\midrule
2&0&\text{linear}&180^\circ\\
3&0&\text{trigonal planar}&120^\circ\\
4&0&\text{tetrahedral}&109.5^\circ\\
4&1&\text{trigonal pyramidal}&\approx107^\circ\\
4&2&\text{bent}&\approx105^\circ\\
5&0&\text{trigonal bipyramidal}&90^\circ,120^\circ\\
6&0&\text{octahedral}&90^\circ\\
6&1&\text{square pyramidal}&\approx90^\circ\\
6&2&\text{square planar}&90^\circ\\
\bottomrule
\end{array}
\]
Exam wording often asks ``planar'', ``non-planar'', ``bond angle'', ``pyramidal'', or ``bent''.

\subsubsection{Planar molecule recognition}
\[
\text{linear, trigonal planar, square planar, and many aromatic rings are planar}
\]
Tetrahedral and trigonal pyramidal centers are non-planar around that atom. A molecule with all atoms in one plane is planar.

\subsubsection{Molecular polarity}
\[
\vec{\mu}_{\mathrm{mol}}=\sum_i\vec{\mu}_i
\]
Polar bonds can cancel in symmetric geometries. Lone pairs often prevent cancellation and make the molecule polar.

\subsection{Exam recipe: molecular geometry, bond angle, polarity}
\textbf{Use when:} the problem asks for shape, bond angle, planar/non-planar, or dipole.

\textbf{Steps:}
\begin{enumerate}
    \item Draw the Lewis structure and count lone pairs on the central atom.
    \item Compute steric number.
    \item Select VSEPR geometry from the table.
    \item Adjust ideal angle downward if lone pairs are present.
    \item For polarity, add bond dipoles by symmetry.
\end{enumerate}

\textbf{Common traps:} counting a double bond as two VSEPR domains; deciding polarity from bond polarity alone; ignoring lone pairs.

\subsection{Bonding formulas}

\subsubsection{Bond polarity from electronegativity}
\[
\Delta\chi=|\chi_A-\chi_B|
\]
Larger electronegativity difference gives a more polar bond. Geometry decides whether the whole molecule is polar.

\subsubsection{Molecular orbital bond order}
\[
\text{bond order}=\frac{N_{\mathrm{bonding}}-N_{\mathrm{antibonding}}}{2}
\]
Higher bond order usually means shorter, stronger bonds. Zero bond order predicts no stable bond in simple MO treatment.
